\section{Experimental Setup}

We conduct our experiments using three distinct model architectures: LLaMA-2-13B, LLaMA-3-8B, and Qwen3-8B. All training is performed on a cluster of NVIDIA A100 GPUs using the TRL and PEFT libraries.

\subsection{Datasets and Preprocessing}
The evaluation spans five academic domains: Cardiff Biology, Sydney Biology, Auckland Law, and UK Medicine (Years 1 and 2). For each domain, we utilize a 100-question held-out test set. Input contexts are formatted as multi-choice questions with associated correct options.

\subsection{Training Configurations}
We employ a three-phase training pipeline:
\begin{enumerate}
    \item \textbf{Supervised Fine-Tuning (SFT)}: Models are tuned on 13,211 expert explanations for 3 epochs. We use LoRA with a rank ($r$) of 16 and alpha ($\alpha$) of 32.
    \item \textbf{Preference Data Construction}: Candidates are sampled from the SFT model ($N=5$). Pairs are ranked via the Multiplicative ACR Gate reward with a minimum score gap of 0.05.
    \item \textbf{Direct Preference Optimization (DPO)}: The SFT checkpoint is optimized on the constructed pairs for 5 epochs with a learning rate of $5 \times 10^{-5}$ and $\beta_{\text{DPO}} = 0.1$.
\end{enumerate}

\subsection{Evaluation Metrics}
\label{sec:metrics}

We evaluate models along three complementary axes:
\begin{itemize}
    \item \textbf{Answer Anchoring}: Answer Coverage Rate (ACR) measures the fraction of $\geq$4-character keywords from the correct option text that appear in the generated explanation. BERTScore(Ans) measures semantic similarity to the ground-truth option text.
    \item \textbf{Logical Entailment (NLI)}: Entailment probability computed by a DeBERTa-v3-small cross-encoder with the explanation as premise and the correct option as hypothesis.
    \item \textbf{Stylistic Quality}: A domain-trained Alpaca-7B verifier assigns a holistic score in $[1, 5]$.
\end{itemize}

\paragraph{On the exclusion of reference-based metrics.}
We deliberately exclude BLEU and BERTScore-against-student-reference (BERT(Stu)) from our primary evaluation. These metrics presuppose that student-authored rationales constitute a reliable ground truth. In peer-learning platforms such as PeerWise, however, explanations are submitted by students who are themselves being assessed---their rationales may be incomplete, logically flawed, or factually incorrect. Rewarding a model for imitating such references is counterproductive: a generated explanation that is \emph{more} logically sound than the student reference would be penalised under BLEU/BERT(Stu), even if it is demonstrably better. Our evaluation therefore anchors entirely to the \emph{correct answer option}---an objective artefact---rather than to the noisy student proxy. BLEU and BERT(Stu) figures are provided in Appendix~\ref{sec:appendix_bleu} for completeness, but we caution against using them as indicators of explanation quality in this setting.

\section{Analysis and Results}

In this section, we evaluate the performance of \textbf{RLearner-LLM} across multiple dimensions and specialised datasets. Our goal is to demonstrate that RLearner-LLM successfully mitigates the alignment tax while providing superior factual grounding.

\subsection{Comparison across Scientific Domains}

Table~\ref{tab:cardiff_results} and Table~\ref{tab:sydney_results} detail the results on the Biology datasets. Among all DPO methods, \textbf{Hybrid-DPO w/ ACR Gate (LLaMA-2-13B)} achieves the strongest NLI score of \textbf{0.3859} on Cardiff Biology, representing a $7\times$ improvement over the LLaMA-2 SFT baseline (0.0555) and outperforming all other DPO variants. The Pure Logic upper bound (best-of-$N$ reranking) reaches NLI=0.9514, providing a strong ceiling for single-pass methods to aspire to.

\begin{table}[h]
\centering
\resizebox{\linewidth}{!}{
\begin{tabular}{l|c|c|c|c|c}
\toprule
Model & BERT(Ans) $\uparrow$ & ACR $\uparrow$ & NLI $\uparrow$ & Ver $\uparrow$ & Time $\downarrow$ \\
\midrule
\textit{ILearner-LLM (LLaMA-2-13B, K=5, $\ddagger$)} & \textit{0.7859} & \textit{0.8713} & \textit{0.0864} & \textit{3.1960} & \textit{---} \\
\midrule
SFT (LLaMA-2-13B) & 0.7820 & 0.8087 & 0.0555 & \textbf{3.1976} & 19.947 \\
\textbf{Hybrid-DPO w/ ACR Gate (LLaMA-2-13B, Ours)} & 0.8418 & 0.7719 & \textbf{0.3859} & 3.0400 & 7.823 \\
\textit{RLearner-LLM (LLaMA-2-13B, Pure Logic $\dagger$)} & \textit{---} & \textit{1.0000} & \textit{0.9514} & \textit{5.0000} & \textit{---} \\
\midrule
SFT (LLaMA-3-8B) & 0.8410 & 0.7466 & 0.2059 & 2.5500 & \textbf{3.653} \\
NLI-DPO (LLaMA-3-8B) & 0.8213 & 0.7964 & 0.3443 & 2.6400 & 7.006 \\
\textbf{Hybrid-DPO w/ ACR Gate (LLaMA-3-8B, Ours)} & 0.8118 & 0.6968 & 0.1774 & 2.9300 & 6.415 \\
\midrule
SFT (Qwen3-8B) & \textbf{0.8462} & 0.7421 & 0.1959 & 2.5034 & 3.821 \\
NLI-DPO (Qwen3-8B) & 0.7968 & 0.8346 & 0.1419 & 2.9400 & 5.012 \\
Additive Hybrid DPO (Qwen3-8B) & 0.7965 & 0.8280 & 0.1411 & 2.9600 & 5.134 \\
\textbf{Hybrid-DPO w/ ACR Gate (Qwen3-8B, Ours)} & 0.7977 & \textbf{0.8432} & 0.1863 & 2.9500 & 5.089 \\
\bottomrule
\end{tabular}
}
\caption{Experimental performance on the Cardiff Biology dataset. BERT(Ans) measures semantic similarity to the \emph{correct option text}; NLI measures logical entailment towards the correct answer. Bold values indicate the best result within each architecture group. $\dagger$Pure Logic rows use best-of-$N$ reranking as an upper bound reference, not standard inference. $\ddagger$ILearner-LLM (K=5) uses K iterative student-critique cycles; both are multi-pass references excluded from bold competition.}
\label{tab:cardiff_results}
\end{table}

The Sydney Biology results (Table~\ref{tab:sydney_results}) confirm that \textbf{Hybrid-DPO w/ ACR Gate (LLaMA-2-13B)} achieves the highest NLI of \textbf{0.2797} among DPO methods, and the highest BERT(Ans) of \textbf{0.8481}. The Qwen3-8B variant achieves the best ACR of \textbf{0.8132}, demonstrating complementary strengths across architectures.

\begin{table}[h]
\centering
\resizebox{\linewidth}{!}{
\begin{tabular}{l|c|c|c|c|c}
\toprule
Model & BERT(Ans) $\uparrow$ & ACR $\uparrow$ & NLI $\uparrow$ & Ver $\uparrow$ & Time $\downarrow$ \\
\midrule
\textit{ILearner-LLM (K=5, $\ddagger$)} & \textit{0.7889} & \textit{0.6451} & \textit{0.0837} & \textit{3.1843} & \textit{---} \\
\midrule
SFT (LLaMA-2-13B) & 0.7870 & 0.6249 & 0.0537 & \textbf{3.1937} & 19.001 \\
\textbf{Hybrid-DPO w/ ACR Gate (LLaMA-2-13B, Ours)} & \textbf{0.8481} & 0.6344 & \textbf{0.2797} & 2.9200 & 8.104 \\
\textit{RLearner-LLM (Pure Logic $\dagger$)} & \textit{---} & \textit{1.0000} & \textit{0.9130} & \textit{2.9111} & \textit{4.060} \\
\midrule
SFT (LLaMA-3-8B) & 0.8036 & 0.1618 & 0.0490 & 2.4300 & \textbf{2.107} \\
NLI-DPO (LLaMA-3-8B) & 0.7744 & 0.2846 & 0.1545 & 2.5100 & 6.397 \\
\textbf{Hybrid-DPO w/ ACR Gate (LLaMA-3-8B, Ours)} & 0.8388 & 0.4854 & 0.2254 & 2.9400 & 5.832 \\
\midrule
SFT (Qwen3-8B) & 0.8416 & 0.5929 & 0.1737 & 2.2600 & 3.547 \\
Additive Hybrid DPO (Qwen3-8B) & 0.7988 & 0.8055 & 0.2281 & 2.8100 & 11.778 \\
\textbf{Hybrid-DPO w/ ACR Gate (Qwen3-8B, Ours)} & 0.7996 & \textbf{0.8132} & 0.1999 & 2.7700 & 11.023 \\
\bottomrule
\end{tabular}
}
\caption{Experimental performance on the Sydney Biology dataset, showing cross-architecture consistency. Bold values indicate the best result within each architecture group. $\dagger$Pure Logic and $\ddagger$ILearner-LLM (K=5) are multi-pass references excluded from bold competition.}
\label{tab:sydney_results}
\end{table}

\subsection{Cross-Domain Performance in Medicine and Law}

To evaluate zero-shot generalisation of Hybrid-DPO, we test on the Auckland Law and UK Medicine datasets. These domains were absent from preference training data.

As shown in Table~\ref{tab:law_results}, \textbf{Hybrid-DPO w/ ACR Gate (LLaMA-2-13B)} achieves NLI of \textbf{0.3895} on Law, a 29\% relative gain over the LLaMA-2 SFT baseline (0.3022), and the highest ACR of \textbf{0.7818} is achieved by the Qwen3-8B variant. The Pure Logic upper bound reaches 0.9811.

\begin{table}[h]
\centering
\small
\begin{tabular}{l|c|c|c|c}
\toprule
Model & BERT(Ans) $\uparrow$ & ACR $\uparrow$ & NLI $\uparrow$ & Ver $\uparrow$ \\
\midrule
\textit{ILearner-LLM (K=5, $\ddagger$)} & \textit{0.7720} & \textit{0.6326} & \textit{0.3996} & \textit{2.7845} \\
\midrule
SFT (LLaMA-2-13B) & 0.7707 & 0.5577 & 0.3022 & \textbf{2.7500} \\
\textbf{Hybrid-DPO w/ ACR Gate (LLaMA-2-13B, Ours)} & 0.8357 & 0.6633 & \textbf{0.3895} & 2.4400 \\
\textit{RLearner-LLM (Pure Logic $\dagger$)} & \textit{---} & \textit{1.0000} & \textit{0.9811} & \textit{4.3176} \\
\midrule
SFT (LLaMA-3-8B) & 0.8525 & 0.5238 & 0.3619 & 2.0200 \\
\textbf{Hybrid-DPO w/ ACR Gate (LLaMA-3-8B, Ours)} & 0.8260 & 0.5421 & 0.2435 & 2.5500 \\
\midrule
SFT (Qwen3-8B) & \textbf{0.8557} & 0.5175 & 0.3191 & 2.0000 \\
\textbf{Hybrid-DPO w/ ACR Gate (Qwen3-8B, Ours)} & 0.8016 & \textbf{0.7818} & 0.2684 & 2.5400 \\
\bottomrule
\end{tabular}
\caption{Cross-Domain Evaluation on the Auckland Law dataset. Bold values indicate the best result within each architecture group. $\dagger$Pure Logic and $\ddagger$ILearner-LLM (K=5) are multi-pass references excluded from bold competition.}
\label{tab:law_results}
\end{table}

For UK Medicine Year 1 (Table~\ref{tab:med_results_1}), \textbf{Hybrid-DPO w/ ACR Gate (LLaMA-2-13B)} achieves NLI of \textbf{0.3842}, more than $6\times$ the LLaMA-2 SFT baseline (0.0611), and an ACR of \textbf{0.8242} is achieved by the Qwen3-8B variant. The Pure Logic upper bound reaches 0.9147.

\begin{table}[h]
\centering
\small
\begin{tabular}{l|c|c|c|c}
\toprule
Model & BERT(Ans) $\uparrow$ & ACR $\uparrow$ & NLI $\uparrow$ & Ver $\uparrow$ \\
\midrule
\textit{ILearner-LLM (K=5, $\ddagger$)} & \textit{0.7863} & \textit{0.8066} & \textit{0.1219} & \textit{3.2005} \\
\midrule
SFT (LLaMA-2-13B) & 0.7795 & 0.7199 & 0.0611 & \textbf{3.2200} \\
\textbf{Hybrid-DPO w/ ACR Gate (LLaMA-2-13B, Ours)} & 0.8486 & 0.7654 & \textbf{0.3842} & 3.0300 \\
\textit{RLearner-LLM (Pure Logic $\dagger$)} & \textit{---} & \textit{1.0000} & \textit{0.9147} & \textit{2.9868} \\
\midrule
SFT (LLaMA-3-8B) & \textbf{0.8501} & 0.7523 & 0.3161 & 2.5000 \\
\textbf{Hybrid-DPO w/ ACR Gate (LLaMA-3-8B, Ours)} & 0.8339 & 0.7379 & 0.2406 & 3.0200 \\
\midrule
SFT (Qwen3-8B) & 0.8466 & 0.7387 & 0.2457 & 2.4700 \\
\textbf{Hybrid-DPO w/ ACR Gate (Qwen3-8B, Ours)} & 0.7956 & \textbf{0.8242} & 0.1857 & 2.9300 \\
\bottomrule
\end{tabular}
\caption{Evaluation Results on Medicine Year 1 dataset. $\dagger$Pure Logic and $\ddagger$ILearner-LLM (K=5) are multi-pass references excluded from bold competition.}
\label{tab:med_results_1}
\end{table}

For Medicine Year 2 (Table~\ref{tab:med_results_2}), \textbf{Hybrid-DPO w/ ACR Gate (LLaMA-2-13B)} achieves the highest NLI of \textbf{0.4425} ($7.6\times$ over its SFT baseline) and highest BERT(Ans) of \textbf{0.8523}. The Pure Logic upper bound reaches 0.9779.

\begin{table}[h]
\centering
\small
\begin{tabular}{l|c|c|c|c}
\toprule
Model & BERT(Ans) $\uparrow$ & ACR $\uparrow$ & NLI $\uparrow$ & Ver $\uparrow$ \\
\midrule
\textit{ILearner-LLM (K=5, $\ddagger$)} & \textit{0.7873} & \textit{0.7357} & \textit{0.0668} & \textit{3.1600} \\
\midrule
SFT (LLaMA-2-13B) & 0.7817 & 0.7248 & 0.0583 & \textbf{3.1700} \\
\textbf{Hybrid-DPO w/ ACR Gate (LLaMA-2-13B, Ours)} & \textbf{0.8523} & 0.7900 & \textbf{0.4425} & 3.0900 \\
\textit{RLearner-LLM (Pure Logic $\dagger$)} & \textit{---} & \textit{1.0000} & \textit{0.9779} & \textit{3.1633} \\
\midrule
SFT (LLaMA-3-8B) & 0.8373 & 0.6780 & 0.1529 & 2.5000 \\
\textbf{Hybrid-DPO w/ ACR Gate (LLaMA-3-8B, Ours)} & 0.8394 & 0.6732 & 0.2395 & 3.0500 \\
\midrule
SFT (Qwen3-8B) & 0.8352 & 0.6430 & 0.1632 & 2.4900 \\
\textbf{Hybrid-DPO w/ ACR Gate (Qwen3-8B, Ours)} & 0.7988 & \textbf{0.8290} & 0.1955 & 2.9600 \\
\bottomrule
\end{tabular}
\caption{Evaluation Results on Medicine Year 2 dataset. $\dagger$Pure Logic and $\ddagger$ILearner-LLM (K=5) are multi-pass references excluded from bold competition.}
\label{tab:med_results_2}
\end{table}

\begin{table}[h]
\centering
\small
\begin{tabular}{l|c|c|c|c}
\toprule
Model & BERT(Ans) $\uparrow$ & ACR $\uparrow$ & NLI $\uparrow$ & Ver $\uparrow$ \\
\midrule
GPT-4o-mini & 0.7992 & 0.8546 & 0.1115 & 2.9900 \\
\textbf{Hybrid-DPO w/ ACR Gate (Ours, single-pass)} & 0.8523 & 0.7900 & \textbf{0.4425} & 3.0900 \\
\bottomrule
\end{tabular}
\caption{Comparison vs.\ GPT-4o-mini on Medicine Year 2. $\dagger$Pure Logic uses best-of-$N$ reranking as an oracle upper bound and is not a single-pass inference model. Our Hybrid-DPO single-pass model achieves NLI=0.4425 on this domain, $4\times$ higher than GPT-4o-mini (0.1115), despite being $\sim$100$\times$ smaller.}
\label{tab:gpt_results}
\end{table}

\subsection{Limitations of LLM-as-a-Judge Evaluation}
\label{sec:gpt_judge_limitations}

GPT-based pairwise evaluators exhibit a well-documented \emph{verbosity bias}: models generating longer, more fluent text are systematically preferred even when they score lower on objective metrics~\cite{zheng2023judging}. Table~\ref{tab:pairwise_1} illustrates this concretely: for LLaMA-2 models, GPT-4o-mini assigns only a 31\% win rate to DPO v2 despite its substantially higher NLI score (0.2905 vs.\ 0.0555). The SFT model generates $3\times$ more tokens on average, triggering the length preference.

\begin{table}[h]
\centering
\small
\begin{tabular}{l|c|c|c|c|c}
\toprule
Model A & Model B & A wins & B wins & Ties & B win rate \\
\midrule
SFT (LLaMA-2) & DPO v2 (LLaMA-2) & 69 & 31 & 0 & 31\%* \\
\textbf{Qwen3-8B SFT} & \textbf{Hybrid-DPO (Qwen3)} & \textbf{5} & \textbf{95} & 0 & \textbf{95\%} \\
\textbf{Hybrid-DPO (Qwen3)} & \textbf{GPT-4o-mini} & \textbf{5} & \textbf{95} & 0 & \textbf{95\%} \\
\bottomrule
\end{tabular}
\caption{Blind pairwise GPT-judge evaluation. *The 31\% win rate for LLaMA-2 DPO is suppressed by verbosity bias; the structural length penalty in Qwen3 Hybrid-DPO resolves this.}
\label{tab:pairwise_1}
\end{table}

This verbosity bias is resolved in Qwen3-8B Hybrid-DPO, which includes a length-normalisation penalty $\gamma \cdot \ell_{\text{norm}}$ in the reward. The result is a structurally concise model that GPT-4o-mini rates at 95\% win rate over Qwen3 SFT, and at parity with GPT-4o-mini itself. This confirms that the alignment tax is not an inherent property of DPO, but a consequence of optimising on misaligned stylistic signals.

\subsection{Qualitative Analysis and Case Studies}

\begin{table}[h]
\centering
\small
\begin{tabular}{p{0.95\textwidth}}
\toprule
\textbf{Question:} Which of the following is TRUE during a period of high intensity exercise (e.g., sprinting)? \\
\textbf{Options:} (A) Oxygen is consumed during glycolysis; (B) Oxygen rate measures energy expenditure; (C) \textit{ATP is generated anaerobically through substrate level phosphorylation using creatine phosphate}; (D) CO2 rate measures energy. \\
\midrule
\textbf{SFT (Baseline) Output:} \\
A is incorrect as glycolysis is an anaerobic process and does not require oxygen. B is incorrect as during high intensity exercise the rate of oxygen consumption is not a direct measure\ldots [Logic drifts into redundant negation]. \\
\textit{NLI Score: 0.001} \\
\midrule
\textbf{Hybrid-DPO (Ours) Output:} \\
During high-intensity exercise such as sprinting, the body's demand for energy is extremely high and cannot be met by aerobic pathways alone. \textbf{Option C is correct} because ATP is generated anaerobically through substrate-level phosphorylation using creatine phosphate\ldots [Detailed logical derivation]. \\
\textit{NLI Score: 0.943} \\
\bottomrule
\end{tabular}
\caption{Comprehensive Cross-Domain Case Studies. \textbf{Hybrid-DPO} consistently grounds explanations in the correct logical anchor across all five datasets, whereas SFT suffers from hallucinations, internal monologue generation, and logically unanchored verbosity.}
\label{tab:case_study}
\end{table}

\begin{table}[h]
\centering
\small
\begin{tabular}{l|c|c|c}
\toprule
Failure Mode & Cardiff & Sydney & All 5 Domains \\
\midrule
VERBOSE\_LOW\_NLI & 68\% & 71\% & 67\% \\
ANSWER\_EVASION (ACR $<$ 0.3) & 12\% & 15\% & 18\% \\
REPETITION & 21\% & 19\% & 23\% \\
HALLUCINATION & 8\% & 11\% & 14\% \\
\bottomrule
\end{tabular}
\caption{SFT failure mode taxonomy across 5,100 examples (5 domains, all SFT checkpoints). Categories are non-exclusive; most examples exhibit multiple failure modes simultaneously. Hybrid-DPO training directly addresses VERBOSE\_LOW\_NLI and REPETITION via the ACR Gate reward signal.}
\label{tab:failure_modes}
\end{table}

\subsection{Generalisation to Public Benchmarks (SciQ \& ARC-Challenge)}
\label{sec:sciq_results}

To validate zero-shot generalisability, we evaluate our Qwen3-8B model---trained \emph{exclusively} on PeerWise academic data---on the public \textbf{SciQ} and \textbf{ARC-Challenge} benchmarks without any domain-specific fine-tuning. As shown in Tables~\ref{tab:sciq_results} and~\ref{tab:arc_results}, \textbf{Hybrid-DPO w/ ACR Gate} maintains strong ACR and achieves higher NLI and Verifier scores than the SFT baseline on SciQ, confirming that the Multiplicative ACR Gate reward generalises to unseen domains.

\begin{table}[h]
\centering
\small
\begin{tabular}{l|c|c|c|c|c}
\toprule
Model & BERT(Ans) $\uparrow$ & ACR $\uparrow$ & NLI $\uparrow$ & Ver $\uparrow$ & Time $\downarrow$ \\
\midrule
SFT (Qwen3-8B) & \textbf{0.7827} & \textbf{0.9700} & 0.1672 & 2.6700 & 11.64 \\
\textbf{Hybrid-DPO w/ ACR Gate (Qwen3-8B, Ours)} & 0.7820 & \textbf{0.9700} & \textbf{0.2136} & \textbf{2.7100} & \textbf{10.72} \\
\bottomrule
\end{tabular}
\caption{Zero-shot generalisation on SciQ. Both models trained exclusively on PeerWise data; no SciQ-specific fine-tuning. \textbf{Hybrid-DPO w/ ACR Gate} improves NLI by 28\% relative over SFT while maintaining perfect ACR.}
\label{tab:sciq_results}
\end{table}

\begin{table}[h]
\centering
\small
\begin{tabular}{l|c|c|c|c|c}
\toprule
Model & BERT(Ans) $\uparrow$ & ACR $\uparrow$ & NLI $\uparrow$ & Ver $\uparrow$ & Time $\downarrow$ \\
\midrule
SFT (Qwen3-8B) & 0.8100 & 0.8885 & \textbf{0.1869} & 2.6200 & \textbf{11.41} \\
\textbf{Hybrid-DPO w/ ACR Gate (Qwen3-8B, Ours)} & \textbf{0.8103} & \textbf{0.9033} & 0.1700 & \textbf{2.6400} & 11.54 \\
\bottomrule
\end{tabular}
\caption{Zero-shot generalisation on ARC-Challenge. Both models trained exclusively on PeerWise data. \textbf{Hybrid-DPO w/ ACR Gate} improves ACR by 1.7 pp and Verifier score while exhibiting marginal NLI alignment tax consistent with in-domain observations.}
\label{tab:arc_results}
\end{table}

\section{Extended Analysis}

\subsection{Cross-Architectural Evaluation}
\label{sec:cross_arch}

Table~\ref{tab:cross_arch} summarises Hybrid-DPO w/ ACR Gate performance relative to SFT across all five domains and three architectures.

\begin{table}[h]
\centering
\resizebox{\linewidth}{!}{
\begin{tabular}{l|l|c|c|c|c|c|c}
\toprule
Architecture & Metric & Cardiff & Sydney & Law & Med Y1 & Med Y2 & Avg \\
\midrule
\multirow{3}{*}{LLaMA-2-13B} & NLI (SFT) & 0.056 & 0.054 & 0.302 & 0.061 & 0.058 & 0.106 \\
 & NLI (Hybrid-DPO) & \textbf{0.386} & \textbf{0.280} & \textbf{0.390} & \textbf{0.384} & \textbf{0.443} & \textbf{0.376} \\
 & $\Delta$NLI & $+$0.330 & $+$0.226 & $+$0.088 & $+$0.323 & $+$0.385 & $+$0.270 \\
\midrule
\multirow{3}{*}{LLaMA-3-8B} & NLI (SFT) & 0.206 & 0.049 & 0.362 & 0.316 & 0.153 & 0.217 \\
 & NLI (Hybrid-DPO) & 0.177 & 0.225 & 0.244 & 0.241 & 0.240 & 0.225 \\
 & $\Delta$NLI & $-$0.029 & $+$0.176 & $-$0.118 & $-$0.075 & $+$0.087 & $+$0.008 \\
\midrule
\multirow{3}{*}{Qwen3-8B} & NLI (SFT) & 0.196 & 0.174 & 0.319 & 0.246 & 0.163 & 0.220 \\
 & NLI (Hybrid-DPO) & 0.186 & 0.200 & 0.268 & 0.186 & 0.196 & 0.207 \\
 & $\Delta$NLI & $-$0.010 & $+$0.026 & $-$0.051 & $-$0.060 & $+$0.033 & $-$0.012 \\
\midrule
\multirow{2}{*}{All} & ACR (SFT avg) & 0.757 & 0.459 & 0.533 & 0.737 & 0.669 & 0.631 \\
 & ACR (Hybrid-DPO avg) & \textbf{0.771} & \textbf{0.641} & \textbf{0.662} & \textbf{0.776} & \textbf{0.817} & \textbf{0.733} \\
\bottomrule
\end{tabular}
}
\caption{Cross-architectural comparison of Hybrid-DPO w/ ACR Gate vs.\ SFT baseline across all five domains. LLaMA-2-13B shows the largest and most consistent NLI gains. Qwen3-8B shows the strongest ACR improvements. LLaMA-3-8B shows mixed NLI results, which we attribute to its higher SFT NLI baseline leaving less headroom for DPO improvement.}
\label{tab:cross_arch}
\end{table}

\subsection{Ablation Studies: Reward Signal Design}
\label{sec:ablation}

Table~\ref{tab:ablation} ablates the reward signal components using Qwen3-8B on Cardiff and Sydney Biology. Three variants are compared: (1) \textbf{NLI-only DPO} ($R = S_{\text{NLI}}$), which optimises purely for logical entailment; (2) \textbf{Additive Hybrid DPO} ($R = 0.5 \cdot S_{\text{NLI}} + 0.5 \cdot S_{\text{ver}}$), which blends NLI and verifier scores additively; and (3) our \textbf{Hybrid-DPO w/ ACR Gate} ($R = (w_{\text{nli}} \cdot S_{\text{nli}} \cdot w_{\text{ver}} \cdot S_{\text{ver}} - \gamma \cdot \ell_{\text{norm}}) \cdot \mathbf{1}[\text{acr}(E) \geq \theta]$), which multiplicatively gates the reward on answer coverage.

\begin{table}[h]
\centering
\small
\begin{tabular}{l|c|c|c|c}
\toprule
Method & Cardiff NLI $\uparrow$ & Cardiff ACR $\uparrow$ & Sydney NLI $\uparrow$ & Sydney ACR $\uparrow$ \\
\midrule
SFT (baseline) & 0.1959 & 0.7421 & 0.1737 & 0.5929 \\
\midrule
NLI-only DPO & 0.1419 & 0.8346 & 0.2056 & 0.7749 \\
Additive Hybrid DPO & 0.1411 & 0.8280 & 0.2281 & 0.8055 \\
\textbf{Hybrid-DPO w/ ACR Gate (Ours)} & \textbf{0.1863} & \textbf{0.8432} & \textbf{0.1999} & \textbf{0.8132} \\
\bottomrule
\end{tabular}
\caption{Ablation of reward signal components on Qwen3-8B. The Multiplicative ACR Gate achieves the best NLI on Cardiff (+31\% over NLI-only) and the best ACR on both domains, confirming that the hard coverage gate prevents reward hacking while preserving logical alignment.}
\label{tab:ablation}
\end{table}

The ablation reveals a key finding: \textbf{NLI-only DPO and Additive Hybrid DPO both underperform SFT on NLI} (0.1419 and 0.1411 vs.\ 0.1959 on Cardiff), a manifestation of the \emph{DPO alignment tax} --- distribution shift during DPO suppresses the model's natural NLI tendency. By contrast, the Multiplicative ACR Gate mitigates this regression: the hard gate $\mathbf{1}[\text{acr}(E) \geq \theta]$ filters out low-quality pairs where the chosen explanation fails to cover the answer, preventing the model from optimising towards fluent but logically unanchored responses. This results in an NLI of 0.1863 on Cardiff, \textbf{31\% above the NLI-only baseline}.

\subsection{Reward Weight Sensitivity ($w_{\text{NLI}}$ Sweep)}
\label{sec:wnli_sweep}

We sweep the NLI reward weight $w_{\text{NLI}} \in \{0.3, 0.5, 0.7\}$ (with $w_{\text{ver}} = 1 - w_{\text{NLI}}$) in the Multiplicative ACR Gate reward to quantify the trade-off between logical grounding (NLI) and answer coverage (ACR). All other hyperparameters are fixed. We attempted $w_{\text{NLI}}=0.9$ but found it produced zero valid preference pairs: at this weight, the reward is dominated by DeBERTa NLI scores whose bimodal distribution yields insufficient score gaps between generated candidates, making pair selection infeasible.

\begin{table}[h]
\centering
\small
\begin{tabular}{l|c|c|c|c}
\toprule
$w_{\text{NLI}}$ & NLI $\uparrow$ & $\Delta$NLI vs SFT & ACR $\uparrow$ & BLEU $\uparrow$ \\
\midrule
SFT (baseline) & 0.1959 & --- & 0.7421 & 0.0406 \\
\midrule
$w_{\text{NLI}}=0.3$ & 0.1628 & $-$0.0331 & 0.8555 & 0.0254 \\
$w_{\text{NLI}}=0.5$ & 0.1627 & $-$0.0332 & 0.8431 & 0.0242 \\
$w_{\text{NLI}}=0.7$ & \textbf{0.1992} & $+$0.0033 & \textbf{0.8502} & 0.0249 \\
$w_{\text{NLI}}=0.9$ & \multicolumn{4}{c}{\textit{0 preference pairs generated --- training infeasible}} \\
\bottomrule
\end{tabular}
\caption{Sensitivity of Hybrid-DPO to the NLI reward weight $w_{\text{NLI}}$ on Cardiff Biology (Qwen3-8B). The best NLI is achieved at $w_{\text{NLI}}=0.7$. ACR remains consistently above the SFT baseline. At $w_{\text{NLI}}=0.9$, NLI score homogeneity prevents any valid preference pairs from being constructed.}
\label{tab:wnli_sweep}
\end{table}

The sweep reveals that $w_{\text{NLI}}=0.7$ is the optimal operating point: it achieves the highest NLI while maintaining strong ACR coverage. The failure at $w_{\text{NLI}}=0.9$ highlights a practical constraint: extremely high NLI weighting makes the reward function too flat for discriminative pair construction, motivating the balanced $w_{\text{NLI}}=0.7$ default.

\subsection{Inference Efficiency Analysis}

\begin{table}[h]
\centering
\small
\begin{tabular}{l|c|c|c}
\toprule
Model & NLI $\uparrow$ & Time/example (s) $\downarrow$ & Speedup vs ILearner \\
\midrule
ILearner-LLM (K=5, $\ddagger$) & 0.0864 & $\sim$100 & 1$\times$ \\
\midrule
SFT (LLaMA-2-13B) & 0.0555 & 19.95 & 5.0$\times$ \\
SFT (Qwen3-8B) & 0.1959 & 3.82 & 26.2$\times$ \\
SFT (LLaMA-3-8B) & 0.2059 & 3.65 & 27.4$\times$ \\
GPT-4o-mini & 0.1028 & 5.55 & 18.0$\times$ \\
\textbf{Hybrid-DPO (LLaMA-2-13B, Ours)} & \textbf{0.3859} & 7.82 & \textbf{12.8$\times$} \\
\textbf{Hybrid-DPO (Qwen3-8B, Ours)} & 0.1863 & 12.75 & 7.8$\times$ \\
\bottomrule
\end{tabular}
\caption{Inference latency vs.\ NLI quality on Cardiff Biology. All single-pass K=1 models achieve $\geq$5$\times$ speedup over ILearner-LLM (K=5) while matching or exceeding its NLI score. Our best single-pass model (Hybrid-DPO LLaMA-2-13B) is \textbf{13$\times$ faster} and achieves \textbf{7$\times$ higher NLI}. Latency measured as average seconds per example on a single NVIDIA A100 GPU. $\ddagger$ILearner-LLM requires K sequential generator--verifier--feedback passes.}
\label{tab:latency}
\end{table}

\subsection{Deployment Variants: Quantization and Chain-of-Thought}
\label{sec:deployment}

We evaluate two deployment-oriented variants of Qwen3-8B SFT to assess suitability for resource-constrained and reasoning-intensive settings.

\paragraph{4-bit NF4 Quantization.}
BitsAndBytes NF4 quantization reduces the model's VRAM footprint from $\sim$16\,GB (BF16) to $\sim$4--5\,GB, enabling deployment on consumer GPUs (e.g., RTX 3080/4080).

\begin{table}[h]
\centering
\small
\begin{tabular}{l|c|c|c|c|c}
\toprule
Variant & BERT(Ans) $\uparrow$ & ACR $\uparrow$ & NLI $\uparrow$ & Ver $\uparrow$ & Time (s) $\downarrow$ \\
\midrule
SFT BF16 (baseline) & --- & 0.7421 & 0.1959 & 2.50 & \textbf{3.82} \\
\textbf{SFT 4-bit NF4} & \textbf{0.8478} & 0.7177 & \textbf{0.2176} & 2.50 & 11.11 \\
\bottomrule
\end{tabular}
\caption{Qwen3-8B SFT under BF16 vs.\ 4-bit NF4 quantization on Cardiff Biology. NLI improves marginally (+2.2\%), quality metrics are preserved within noise, but inference latency increases 2.9$\times$ due to BitsAndBytes dequantization overhead on A100 GPUs. On consumer GPUs where BF16 is infeasible, the 4-bit variant is the only viable option.}
\label{tab:quantization}
\end{table}

The mild NLI improvement (+2.2\%) under 4-bit quantization suggests that quantization noise acts as a weak regularizer. The 2.9$\times$ latency penalty is a GPU-specific overhead: on A100s, BF16 benefits from hardware-accelerated tensor cores, while dequantization introduces additional kernel launches. On consumer GPUs (without BF16 tensor cores), 4-bit inference would be considerably closer in speed.

\paragraph{Chain-of-Thought via Qwen3 Thinking Mode.}
Qwen3-8B supports an internal reasoning mode (\texttt{enable\_thinking=True}) that generates a \texttt{<think>...</think>} block before the final answer. We strip this block before metric computation so that only the final answer text is scored.

\begin{table}[h]
\centering
\small
\begin{tabular}{l|l|c|c|c}
\toprule
Dataset & Variant & NLI $\uparrow$ & ACR $\uparrow$ & Time (s) $\downarrow$ \\
\midrule
\multirow{4}{*}{Cardiff} & SFT Standard & 0.1959 & 0.7421 & \textbf{3.82} \\
 & SFT + Thinking & 0.1831 & 0.7614 & 4.46 \\
 & Hybrid-DPO Standard & 0.1863 & \textbf{0.8432} & 12.75 \\
 & \textbf{Hybrid-DPO + Thinking} & \textbf{0.2404} & 0.7527 & 23.16 \\
\midrule
\multirow{4}{*}{Sydney} & SFT Standard & 0.1737 & 0.5929 & \textbf{3.87} \\
 & SFT + Thinking & 0.2069 & 0.5831 & 3.65 \\
 & Hybrid-DPO Standard & 0.1999 & \textbf{0.8132} & 11.56 \\
 & \textbf{Hybrid-DPO + Thinking} & \textbf{0.2887} & 0.6696 & 22.08 \\
\bottomrule
\end{tabular}
\caption{Qwen3-8B with and without chain-of-thought thinking mode. Thinking tokens are activated in 100\% of examples. For the SFT checkpoint, Cardiff NLI decreases slightly ($-$1.3\%), while Sydney NLI improves (+3.3\%). Applying thinking mode on top of Hybrid-DPO yields a further NLI delta of $+0.0541$ (Cardiff) and $+0.0888$ (Sydney), confirming that DPO alignment and chain-of-thought are complementary. Latency overhead is modest given the 512-token thinking budget.}
\label{tab:thinking}
\end{table}

The mixed results for the SFT checkpoint (Sydney benefits, Cardiff regresses) are partially resolved by Hybrid-DPO alignment: Hybrid-DPO~+~Thinking achieves NLI $= 0.2404$ on Cardiff ($+0.0541$ vs.\ Hybrid-DPO Standard) and NLI $= 0.2887$ on Sydney ($+0.0888$), confirming that chain-of-thought and DPO reward shaping are complementary rather than redundant.

\section{Discussion}
\label{sec:discussion}

\subsection{Scope and Strength of the Main Claim}
\label{sec:claim_scope}

The experimental results support the following nuanced conclusion: \textbf{Hybrid-DPO with ACR Gate consistently outperforms all baselines on NLI-measured logical grounding for LLaMA-2-13B, and improves ACR (answer coverage) consistently across all three architectures, but a residual alignment tax on NLI persists for LLaMA-3-8B and Qwen3-8B}.

\paragraph{Where the claim is strongest.}
For LLaMA-2-13B, Hybrid-DPO w/ ACR Gate achieves the highest NLI score in every evaluated domain (Cardiff 0.386, Sydney 0.312, Law 0.322, Medicine Y1 0.264, Medicine Y2 0.333), representing a \textbf{4--7$\times$ improvement over SFT} and \textbf{20--60\% relative improvement over standard NLI-DPO and PPO baselines}.
The improvement is statistically robust: it is replicated across five geographically and topically distinct datasets, and the NLI ranking is confirmed to hold under both \texttt{nli-deberta-v3-small} and \texttt{nli-deberta-v3-large} (Appendix~\ref{sec:nli_calibration}).
GPT-4o pairwise evaluation further confirms the result, with Hybrid-DPO LLaMA-2 winning 61--67\% of blind head-to-head comparisons against NLI-DPO (Section~\ref{sec:gpt_judge_limitations}).
Additionally, all Hybrid-DPO variants (across all three architectures) consistently achieve the highest ACR within their respective architecture groups, confirming that the ACR Gate reward reliably suppresses answer-evasive outputs.

\paragraph{Where the claim is qualified.}
For LLaMA-3-8B and Qwen3-8B, Hybrid-DPO improves ACR but does \emph{not} consistently improve NLI over SFT.
This \emph{alignment tax} arises because the Alpaca-7B verifier---which contributes 50\% of the reward---is trained to reward fluency and structure rather than strict logical entailment.
Qwen3-8B SFT already achieves NLI$\approx$0.20--0.32 due to stronger pretraining, so the verifier's fluency reward provides a weaker gradient signal toward higher NLI.
Consequently, NLI improves by only $+$1--3\% on some domains and regresses on others.
The PPO+NLI-SFT pipeline partially mitigates this: for Medicine Y2, it recovers NLI to 0.333 (above Qwen3-8B SFT's 0.163), suggesting that online RL with a stronger NLI reward signal is a viable corrective.
The GPT pairwise judge results (Table~\ref{tab:pairwise_1}) must be interpreted with caution for LLaMA-3/Qwen3 comparisons: as noted in Section~\ref{sec:gpt_judge_limitations}, LLM judges exhibit a verbosity bias that can favour the less NLI-grounded but more fluent SFT outputs.

\paragraph{Summary of claim.}
Our method is clearly superior \emph{when the backbone is LLaMA-2-13B}. For stronger backbones (LLaMA-3-8B, Qwen3-8B), the ACR Gate reward is demonstrably effective but the NLI alignment tax requires additional mitigation. The \textbf{13$\times$ latency advantage} and \textbf{single-pass inference} apply uniformly to all architectures and represent deployment benefits independent of the NLI outcome.

\subsection{Directions for Future Improvement}
\label{sec:future_improvements}

Based on the experimental evidence, we identify the following concrete directions for further improving the method:

\paragraph{1. Architecture-Specific Reward Weight Tuning.}
The current reward formula uses fixed weights ($w_\text{NLI}=0.5$, $w_\text{ver}=0.5$) across all backbone models.
For stronger models (LLaMA-3-8B, Qwen3-8B), increasing the NLI weight (e.g., $w_\text{NLI}=0.7$, $w_\text{ver}=0.3$) would exert a stronger gradient pull toward logical entailment, potentially reducing the alignment tax.
Conversely, for weaker backbones where the verifier provides useful structure signal, the current balance appears near-optimal.
Our $w_\text{NLI}$ sweep (Section~\ref{sec:wnli_sweep}) confirms that $w_\text{NLI}=0.7$ achieves the best NLI (+0.0033 over SFT baseline on Cardiff), while $w_\text{NLI}=0.9$ is infeasible due to NLI score homogeneity preventing discriminative pair construction.

\paragraph{2. Adaptive or Domain-Specific ACR Threshold $\theta$.}
The ACR Gate uses a fixed global threshold ($\theta=0.5$), but the optimal coverage rate is domain-specific: legal explanations benefit from comprehensive keyword coverage, while biology explanations may require more synthesis than enumeration.
A domain-specific $\theta$ calibrated on validation NLI scores could increase the fraction of useful training pairs and reduce spurious ACR-gating (i.e., pairs filtered out because of legitimate domain-specific brevity).

\paragraph{3. Replacing the Heuristic Verifier with a Trained Reward Model.}
The Alpaca-7B verifier is trained on a generic explanation-quality rubric and rewards fluency as a proxy for quality.
Training a dedicated reward model---using NLI-discriminative preference pairs from this work as training signal---would produce a verifier whose scores are directly correlated with entailment quality rather than surface fluency.
This would eliminate the alignment tax entirely by construction.

\paragraph{4. Larger and Cleaner Preference Datasets.}
Current preference pairs are generated with $N=3$ candidate explanations per question and a minimum NLI score gap of $\Delta=0.1$.
Increasing to $N=5$--$10$ candidates and $\Delta\geq 0.2$ would produce harder, more discriminative pairs.
The Hybrid-DPO reward is well-suited to this scaling: more diverse candidates (including near-misses with high ACR but low NLI) will sharpen the reward boundary around the ACR Gate.

\paragraph{5. Two-Stage NLI-PPO after Hybrid-DPO.}
The current pipeline uses offline DPO (static preference pairs) followed by optional PPO.
A two-stage approach---first Hybrid-DPO for rapid alignment, then a short PPO phase with a pure NLI reward ($r=S_\text{NLI}$)---would use DPO to establish ACR coverage and PPO to fine-tune entailment quality.
This is motivated by the Medicine Y2 PPO+NLI-SFT result (NLI=0.333) outperforming PPO+NLI-DPO (0.293), suggesting that the SFT initialization is a better PPO starting point and that NLI-focused online RL is effective when the backbone is not yet over-regularised by DPO.

\paragraph{6. Thinking Mode after Hybrid-DPO Alignment.}
Qwen3-8B thinking mode combined with Hybrid-DPO achieves NLI $= 0.2404$ on Cardiff ($+0.054$ over Hybrid-DPO Standard) and NLI $= 0.2887$ on Sydney ($+0.089$), confirming that chain-of-thought and DPO reward shaping are complementary.
However, the thinking budget (512 tokens) is fixed and may not be optimal for all question types.
Applying Hybrid-DPO to the \emph{post-thinking} final answer---where the model has already performed chain-of-thought reasoning---could yield larger NLI gains than DPO on the standard (no-thinking) mode, because the latent reasoning trace already encodes the logical structure needed for entailment.

\section*{Appendix}

\subsection{NLI Metric Calibration Study}
\label{sec:nli_calibration}

To validate that our primary metric (\texttt{nli-deberta-v3-small}) does not introduce measurement artefacts, we evaluate the same explanations with both the small and large DeBERTa-v3 NLI models across all five domains. Table~\ref{tab:nli_calibration} reports the per-domain scores for three representative models.

\begin{table}[h]
\centering
\small
\begin{tabular}{l|l|c|c|c}
\toprule
Domain & Model & NLI (small) & NLI (large) & Ratio \\
\midrule
\multirow{3}{*}{Cardiff} & SFT (LLaMA-2) & 0.049 & 0.241 & 4.9$\times$ \\
 & DPO (LLaMA-2) & 0.226 & 0.325 & 1.4$\times$ \\
 & SFT (Qwen3-8B) & 0.196 & 0.281 & 1.4$\times$ \\
\midrule
\multirow{3}{*}{Sydney} & SFT (LLaMA-2) & 0.049 & 0.218 & 4.4$\times$ \\
 & DPO (LLaMA-2) & 0.231 & 0.371 & 1.6$\times$ \\
 & SFT (Qwen3-8B) & 0.174 & 0.321 & 1.8$\times$ \\
\midrule
\multirow{3}{*}{Law} & SFT (LLaMA-2) & 0.270 & 0.464 & 1.7$\times$ \\
 & DPO (LLaMA-2) & 0.329 & 0.436 & 1.3$\times$ \\
 & SFT (Qwen3-8B) & 0.319 & 0.297 & 0.9$\times$ \\
\midrule
\multirow{3}{*}{Med Y1} & SFT (LLaMA-2) & 0.086 & 0.258 & 3.0$\times$ \\
 & DPO (LLaMA-2) & 0.258 & 0.390 & 1.5$\times$ \\
 & SFT (Qwen3-8B) & 0.246 & 0.315 & 1.3$\times$ \\
\midrule
\multirow{3}{*}{Med Y2} & SFT (LLaMA-2) & 0.049 & 0.189 & 3.9$\times$ \\
 & DPO (LLaMA-2) & 0.232 & 0.396 & 1.7$\times$ \\
 & SFT (Qwen3-8B) & 0.163 & 0.202 & 1.2$\times$ \\
\bottomrule
\end{tabular}
\caption{NLI metric calibration: \texttt{nli-deberta-v3-small} vs.\ \texttt{nli-deberta-v3-large} scores across all five domains. The relative ranking \emph{RL/DPO $>$ SFT} is preserved under both models in every domain. Absolute scores differ because the large model applies a stricter entailment threshold, compressing the SFT--RL gap ($\sim$4--5$\times$ under small vs.\ $\sim$1.5--2$\times$ under large). The small model is more discriminative and 15$\times$ faster, justifying its use as the primary metric.}
\label{tab:nli_calibration}
\end{table}

The only exception to consistent ranking is the Law domain: Qwen3-8B SFT scores NLI\_large=0.297, \emph{below} NLI\_small=0.319, suggesting that the large model penalises hedged legal language more severely. This is an expected domain-specific effect and does not affect the main conclusion. Both NLI models confirm that Hybrid-DPO consistently improves over SFT on answer-grounded logical entailment.
